% ID: MAT-GAN-ELL-002
% Titolo: Famiglia di ellissi con parametro
% Disciplina: Matematica
% Area: geometria_analitica
% Argomento: ellisse
% Classe: Terza o quarta secondaria di secondo grado
% Difficolta: 4
% Tipo: problema_numerico
% Tag: matematica, geometria_analitica, ellisse
% Fonte: verifica_ellisse_anonimizzata
% Provenienza: materiale fornito dal docente; testo anonimizzato; soluzione redatta per Exergo

\begin{esercizio}
Considera la seguente equazione:
    \[(k-3)x^2+(k+3)y^2 = k^2-9.\]
    Determina i valori del parametro $k$ tali che:
        \begin{itemize}
            \item l'ellisse ha i fuochi appartenenti all'asse $x$;
            \item l'ellisse ha l'asse maggiore lungo $8$;
            \item l'ellisse ha eccentricità $\dfrac{\sqrt{3}}{2}$;
            \item l'ellisse ha i fuochi appartenenti all'asse $y$.
        \end{itemize}
\end{esercizio}

\begin{soluzione}
Per $k\ne\pm3$ l'equazione diventa $x^2/(k+3)+y^2/(k-3)=1$.
È un'ellisse reale non degenere soltanto per $k>3$.
In questo intervallo $a^2=k+3$, $b^2=k-3$, $c^2=6$, asse maggiore lungo $x$.
Quindi: fuochi sull'asse $x$ per ogni $k>3$;
asse maggiore 8 se $2\sqrt{k+3}=8$, cioè $k=13$;
eccentricità $\sqrt3/2$ se $6/(k+3)=3/4$, cioè $k=5$;
nessun valore dà un'ellisse con fuochi sull'asse $y$.
Per $k=3$ si ha la retta doppia $y=0$, per $k=-3$ la retta doppia $x=0$.
\end{soluzione}
